%====================================================================
%  Stage 2: Endogenous Congestion Spatial Equilibrium (full specification)
%  Beijing gate5 — builds on Stage 1; observed-anchored congestion (v2026-06-29)
%====================================================================
\documentclass[12pt]{article}

\usepackage[margin=1in]{geometry}
\usepackage{amsmath,amssymb,amsthm,mathtools}
\usepackage{graphicx,booktabs,tabularx,array,setspace,float,enumitem}
\usepackage[colorlinks=true,linkcolor=blue,citecolor=blue,urlcolor=blue]{hyperref}
\setstretch{1.15}
\setlength{\parindent}{1.5em}

\newtheorem{definition}{Definition}
\newtheorem{assumption}{Assumption}
\newtheorem{remark}{Remark}

\newcommand{\N}{\mathcal N}
\newcommand{\E}{\mathcal E}
\newcommand{\Lbar}{\bar L}
\newcommand{\Wbar}{\bar W}
\newcommand{\ubar}{\bar u}
\newcommand{\Abar}{\bar A}
\newcommand{\obs}{\mathrm{obs}}
\newcommand{\ff}{\mathrm{ff}}
\newcommand{\corr}{\mathrm{corr}}

\title{Stage~2: Endogenous Congestion\\[0.4em]
\large Beijing Morning Peak, \texttt{model\_input\_edgeproj\_v2\_lcc\_gate5\_v1}\\[0.2em]
\normalsize Observed-anchored congestion, level-form solver, full CF suite}
\author{Yige Duan \and Xiaoruo Feng \and Yizhen Gu}
\date{\today}

\begin{document}
\maketitle

\begin{abstract}
\noindent
This note documents the \textbf{current} Stage~2 implementation: a level-form
Allen--Arkolakis spatial equilibrium on the gate5 Beijing network with
\textbf{link-level congestion feedback} activated.
It extends Stage~1 (\path{FixedDirectedCost_Stage1_gate5.tex}), which held
$\tau_{kl}=\tau^{\obs}_{kl}$ fixed.
Stage~2 keeps the same input bundle, inversion of $(\ubar,\Abar)$ on observed
marginals, and counterfactual taxonomy, but replaces the fixed-cost assumption
with a nested fixed point in which model-implied directed traffic $\Xi_{kl}$
and travel times $\tau_{kl}$ co-evolve.

The \textbf{baseline congestion specification} adopted in the current code is
the \emph{observed-anchored} (exact-hat) update: baseline travel times match
$\tau^{\obs}$ by construction, and the Beijing literature elasticity
$\lambda_{\mathrm{phys}}=1.1$ acts only on \emph{relative} deviations of
$(\Xi,n)$ from the observed allocation.
We document the full equation system, estimation status (external calibration
only), nested algorithm, transmission paths, counterfactual suite, code map,
model-validation diagnostics (implied $\Xi$, OD fit), and tidal-lane
short-run vs.\ long-run congestion-cost decomposition.
\end{abstract}

\tableofcontents
\newpage

%====================================================================
\section{Relation to Stage~1 and scope}\label{sec:scope}
%====================================================================

Stage~1 answered: \emph{under fixed $\tau^{\obs}$, how much do directed vs.\
symmetric calibration worlds differ in implied OD, access, and post-calibration
transport shocks?}
Stage~2 adds the feedback loop missing in Stage~1:
\begin{equation}\label{eq:feedback_loop}
  \underbrace{(l^R,l^F,\Pi,P)}_{\text{sorting \& access}}
  \;\Longrightarrow\;
  \underbrace{\Xi_{kl}}_{\text{(E5) link traffic}}
  \;\Longrightarrow\;
  \underbrace{\tau_{kl}}_{\text{(E6) congestion}}
  \;\Longrightarrow\;
  \underbrace{T,\tau_{ij}}_{\text{route aggregation}}
  \;\Longrightarrow\;
  \text{back to sorting}.
\end{equation}

\paragraph{What carries over unchanged from Stage~1.}
Same $N=1104$ nodes, $|\E|=2642$ directed edges, $\Lbar\approx 88{,}110$ road
commuters, $\tau^{\obs}_{kl}$ from \texttt{edgeproj\_v2} + LCC + detour gate
$G=5$.
Same inversion step: recover $(\ubar_i^{\theta},\Abar_j^{\theta},\chi)$ that
rationalise observed residence and workplace \emph{marginals} at $\tau^{\obs}$.
Same policy taxonomy (uniform speed, hub access, hub productivity, outer amenity,
fringe transport) and tidal-lane pair selection (top-$K$ asymmetric corridors,
$\Delta n$ lane reallocation).

\paragraph{What is new in Stage~2.}
\begin{enumerate}[leftmargin=2em]
  \item \textbf{(E5)--(E6)} endogenous $\Xi_{kl}$ and $\tau_{kl}$.
  \item Two \textbf{congestion modes} for every CF: \emph{exogenous} $\tau$
        (Stage~1 replica on the level-form stack) vs.\ \emph{endogenous} $\tau$.
  \item \textbf{Observed-anchored} congestion as the preferred endogenous spec
        (§\ref{sec:anchored}).
  \item \textbf{Short-run / long-run} horizons for tidal lanes (pinned vs.\ free
        sorting).
  \item \textbf{Congestion-cost decomposition} beyond $\Wbar$
        (§\ref{sec:tidal_diagnostics}).
\end{enumerate}

\paragraph{Headline estimands.}
As in Stage~1, full bilateral OD replication is \emph{not} a success criterion
(marginal fit only).
Under anchored congestion, \textbf{baseline levels} of $\tau$ are pinned;
\textbf{counterfactuals} should be read as \emph{relative changes}:
$\Delta\Wbar$, OD-weighted route-cost changes, treated-edge $\Delta\tau$,
and $(\Xi/\Xi^{\obs})$ responses.

%====================================================================
\section{Environment and notation}\label{sec:data}
%====================================================================

Directed edges $(k,l)\in\E$ carry length $\ell_{kl}$ (km), lane count $n_{kl}$,
free-flow time $\tau^{\ff}_{kl}=\ell_{kl}/v^0_{kl}$, and observed peak time
$\tau^{\obs}_{kl}$.
Iceberg cost $t_{kl}=(\tau_{kl})^{\theta}$ with Frechet shape $\theta=6.83$.

\begin{table}[htbp]
  \centering
  \caption{External parameters (AA Seattle placeholders, gate5)}
  \begin{tabular}{llr}
  \toprule
  Symbol & Meaning & Value \\
  \midrule
  $\theta$ & Route dispersion / iceberg exponent & $6.83$ \\
  $\alpha$ & Workplace spillover elasticity & $-0.12$ \\
  $\beta$  & Residence spillover elasticity & $-0.10$ \\
  $\lambda_{\mathrm{phys}}$ & Beijing inverse-speed elasticity (Barwick et al.) & $\approx 1.1$ \\
  $\lambda_{\mathrm{AA}}$ & Iceberg congestion exponent $=\lambda_{\mathrm{phys}}/\theta$ & $\approx 0.161$ \\
  $\delta_1$ (legacy) & AA Seattle BPR exponent in Stage~1 tidal BPR & $0.488$ \\
  \bottomrule
  \end{tabular}
\end{table}

Input bundle:
\texttt{data\_work/outputs/model\_input\_edgeproj\_v2\_lcc\_gate5\_v1/}.

%====================================================================
\section{Model}\label{sec:model}
%====================================================================

\subsection{Route aggregation and bilateral costs}

Define the adjacency system with uniform rescaling $\kappa$ (spectral guard):
\begin{equation}
  A_{kl} = \frac{1}{\kappa\,\tau_{kl}},
  \qquad
  T = (I-A)^{-1},
  \qquad
  \tau_{ij}^{-\theta} = T_{ij}.
  \label{eq:route_agg}
\end{equation}
On gate5, $\kappa=1$ under anchored baseline.

\subsection{Sorting, market access, and welfare}

Residence and workplace shares $(l_i^R,l_j^F)$ solve the level-form
(E\_R)--(E\_F) system with inverted fundamentals and scale $\chi$ from Step~4:
\begin{align}
  (l_i^R)^{1-\beta\theta} &= \chi\,\ubar_i^{\theta}\sum_j \tau_{ij}^{-\theta}\,
    \Abar_j^{\theta}\,(l_j^F)^{\alpha\theta}, \label{eq:ER}\\
  (l_j^F)^{1-\alpha\theta} &= \chi\,\Abar_j^{\theta}\sum_i \tau_{ij}^{-\theta}\,
    \ubar_i^{\theta}\,(l_i^R)^{\beta\theta}. \label{eq:EF}
\end{align}
Market access:
\begin{align}
  \Pi_i^{-\theta} &= \sum_j \tau_{ij}^{-\theta}\,L_j^F\,P_j^{\theta},
    \label{eq:Pi}\\
  P_j^{-\theta}   &= \sum_i \tau_{ij}^{-\theta}\,L_i^R\,\Pi_i^{\theta}.
    \label{eq:P}
\end{align}
Aggregate welfare (reported, not headline):
\begin{equation}
  \Wbar = \Big[\sum_i \ubar_i^{\theta}\,(L_i^R)^{\beta\theta}\,
    \Pi_i^{-\theta}\Big]^{1/\theta}.
\end{equation}

\subsection{Directed link traffic (E5)}

\begin{equation}
  \Xi_{kl} = t_{kl}^{-\theta}\,P_k^{-\theta}\,\Pi_l^{-\theta},
  \qquad t_{kl} = \tau_{kl}^{1/\theta}.
  \label{eq:E5}
\end{equation}
\emph{Interpretation.} $\Xi_{kl}$ is \textbf{model-assigned expected commuter
mass} on directed link $(k,l)$ implied by Frechet routing and market access.
It is \textbf{not} observed link counts nor vehicles/lane-km.
It correlates with destination employment ($\corr(\ln\Xi,\ln L^F_j)\approx +0.46$)
but \textbf{not} with observed congestion ($\corr(\ln\Xi,\ln(\tau^{\obs}/\tau^{\ff}))\approx +0.08$).

\subsection{Congestion technology (E6): three implemented forms}\label{sec:anchored}

\paragraph{Legacy level form (deprecated for baseline).}
\begin{equation}
  \tau_{kl} \leftarrow \max\Big\{\tau^{\ff}_{kl},\;
    \tau^{\ff}_{kl}\Big(\frac{\Xi_{kl}}{n_{kl}}\Big)^{\lambda}\Big\}.
  \label{eq:E6_legacy}
\end{equation}
Fails to reproduce $\tau^{\obs}$ when $\lambda$ is borrowed from Beijing
literature: raw $\Xi/n$ has no physical density scale.

\paragraph{Density-normalised level form (robustness).}
\begin{equation}
  \tau_{kl} \leftarrow \max\Big\{\tau^{\ff}_{kl},\;
    \tau^{\ff}_{kl}\Big(\frac{\Xi_{kl}}{n_{kl}\ell_{kl}}\Big)^{\lambda_{\mathrm{phys}}}\Big\}.
  \label{eq:E6_density}
\end{equation}
Aligns units with vehicles/lane-km but still does not pin baseline $\tau$.

\paragraph{Observed-anchored form (\textbf{current baseline}).}
Compute observation anchors at fixed $\tau^{\obs}$:
\begin{equation}
  \Xi^{\obs}_{kl} := \Xi_{kl}\big|_{\tau=\tau^{\obs},\,(l^R,l^F,\Pi,P)\,\text{at obs allocation}}.
  \label{eq:Xi_obs}
\end{equation}
Then update:
\begin{equation}
  \boxed{
  \tau_{kl} \leftarrow \max\left\{\tau^{\ff}_{kl},\;
    \tau^{\obs}_{kl}\left(
      \frac{\Xi_{kl}}{\Xi^{\obs}_{kl}}\cdot
      \frac{n^{\obs}_{kl}}{n_{kl}}
    \right)^{\lambda_{\mathrm{phys}}}
  \right\}.
  }
  \label{eq:E6_anchored}
\end{equation}
At the observation point ($\Xi=\Xi^{\obs}$, $n=n^{\obs}$), $\tau=\tau^{\obs}$
by construction.
$\lambda_{\mathrm{phys}}=1.1$ (Barwick et al.\ Beijing) governs \emph{response
to relative load/capacity changes} around that point---consistent with exact-hat
counterfactual logic.

\begin{remark}[What ``anchored'' does and does not fix]
Anchoring fixes \textbf{level misalignment} of $\tau$ and the arbitrary scale of
$\Xi$.
It does \textbf{not} make $\Xi^{\obs}$ a validated traffic count; baseline OD
matrix fit remains poor (§\ref{sec:validation}).
CF credibility rests on \textbf{relative} cost changes, not on $\Xi$ levels
matching road sensors.
\end{remark}

\subsection{Congestion modes in counterfactuals}

\begin{itemize}[leftmargin=2em]
  \item \textbf{Exogenous $\tau$} (\texttt{endogenous\_congestion=False}):
        $\tau_{kl}$ fixed at policy-adjusted values; re-solve
        \eqref{eq:ER}--\eqref{eq:P} only. Replicates Stage~1 shocks on the
        level-form solver.
  \item \textbf{Endogenous $\tau$}: outer loop on \eqref{eq:E6_anchored}
        (or legacy/density variants); transport shocks scale $\tau^{\ff}$;
        lane shocks change $n_{kl}$.
  \item \textbf{Short-run tidal}: pin $(l^R,l^F)$ at baseline while $\tau$
        updates (mechanical + partial GE on network).
  \item \textbf{Long-run tidal}: free sorting after lane shock.
\end{itemize}

%====================================================================
\section{Estimation and calibration}\label{sec:estimation}
%====================================================================

There is \textbf{no structural estimation} of $(\theta,\alpha,\beta,\lambda)$ in
the current pipeline---all are external placeholders pending alignment with
coauthors and Beijing-specific work.

\subsection{Step 1: Primitives (\texttt{build\_primitives.py})}

Load nodes, directed edges, OD flows from
\texttt{model\_input\_edgeproj\_v2\_lcc\_gate5\_v1}.
Export $\ell_{kl},v^0_{kl},v^{\obs}_{kl},\tau^{\ff}_{kl},\tau^{\obs}_{kl},
n_{kl},L^{\obs}_{ij},L^{\obs,R}_i,L^{\obs,F}_j,\Lbar$.

\subsection{Step 2: Parameters (\texttt{parameters.py})}

Frozen dataclass \texttt{Parameters} with AA Seattle $(\theta,\alpha,\beta)$ and
optional \texttt{with\_lambda()} override.
Beijing mapping: use $\lambda_{\mathrm{phys}}=1.1$ in \eqref{eq:E6_anchored}
via \texttt{congestion\_lambda=1.1} in the solver (not $\lambda_{\mathrm{AA}}=0.161$
unless using iceberg-normalised power on $\Xi/n$ without anchoring).

\subsection{Step 3: $\lambda$ diagnostic (\texttt{lambda\_diagnostic.py})}

Cross-section regression (lane-bin FE), diagnostic only:
\begin{equation}
  \ln(\tau^{\obs}_{kl}/\tau^{\ff}_{kl})
  = \lambda\,\ln(\hat\Xi_{kl}/n_{kl}) + \varepsilon_{kl}.
\end{equation}
At legacy placeholder $\lambda=0.071$, $\hat\lambda\approx 0.017$ on gate5---confirms
model $\hat\Xi$ is a weak proxy for observed congestion.
\textbf{Not} used to set $\lambda_{\mathrm{phys}}=1.1$ in anchored spec.

\subsection{Step 4: Inversion (\texttt{invert\_fundamentals.py})}

\begin{assumption}[Observed marginals are equilibrium moments]
\label{ass:marginals}
At $\tau=\tau^{\obs}$, observed shares $(l^{R,\obs},l^{F,\obs})$ equal the
(E\_R)--(E\_F) system for some $(\ubar_i^{\theta},\Abar_j^{\theta},\chi)$.
\end{assumption}
Alternating updates on $(\ubar^{\theta},\Abar^{\theta})$ with geometric-mean
normalisation; output held fixed in all CFs.
Marginal fit: L1 $\lesssim 0.03\%$ on gate5.

\subsection{Step 5: Equilibrium solve (\texttt{equilibrium\_solve.py})}

See §\ref{sec:algorithm}.
No exact-hat algebra: every scenario re-solves levels from $(\ubar,\Abar,\chi)$.

%====================================================================
\section{Nested algorithm}\label{sec:algorithm}
%====================================================================

\subsection{Outer loop (congestion)}

\begin{enumerate}[leftmargin=2em]
  \item Initialise $\tau_{kl}\leftarrow\tau^{\obs}_{kl}$ (or CF-adjusted seed).
  \item \textbf{Inner A}: build $A(\tau)$, $T$, solve \eqref{eq:ER}--\eqref{eq:EF}
        for $(l^R,l^F)$ unless short-run pins apply.
  \item \textbf{Inner B}: solve \eqref{eq:Pi}--\eqref{eq:P} for $(\Pi,P)$.
  \item Compute $\Xi$ from \eqref{eq:E5}.
  \item If exogenous mode: stop.
  \item Else update $\tau$ via \eqref{eq:E6_anchored} (or legacy/density).
  \item \textbf{Damping} in log space:
        $\log\tau \leftarrow (1-d)\log\tau + d\log\tau_{\mathrm{new}}$.
        For $\lambda_{\mathrm{phys}}=1.1$ anchored, use $d=0.25$ (smaller steps
        prevent oscillation; $d=0.5$ may fail to converge).
  \item Repeat until $\max|\log\tau^{(k+1)}-\log\tau^{(k)}|<10^{-5}$.
\end{enumerate}

\subsection{Observation anchors (\texttt{compute\_observation\_anchors})}

One exogenous-$\tau$ pass at $\tau^{\obs}$ yields $(\Xi^{\obs},l^R,l^F,\Pi,P)$
stored for \eqref{eq:E6_anchored}.
Anchors recomputed when fundamentals or $\tau^{\obs}$ change.

\subsection{Counterfactual execution order}

\begin{enumerate}[leftmargin=2em]
  \item Invert fundamentals (once).
  \item Baselines: exogenous $\tau^{\obs}$ and endogenous anchored.
  \item Policy CFs: each policy $\times$ (exog/endog) relative to \emph{own} baseline.
  \item Structural CF1 (geometric symmetric $\tau$ init), CF2 (tidal lanes),
        CF3 (uniform $n$ expansion) --- endogenous anchored.
  \item Tidal grid: endogenous ST/LT; exogenous BPR grid (Stage~1 fixed-$\xi$
        on treated edges only).
\end{enumerate}

%====================================================================
\section{Transmission mechanisms}\label{sec:transmission}
%====================================================================

\subsection{Policy shocks}

\begin{table}[htbp]
  \centering
  \caption{Transmission map (endogenous anchored)}
  \small
  \begin{tabular}{p{3.2cm}p{4.2cm}p{5.5cm}}
  \toprule
  Shock & Direct primitive & GE chain \\
  \midrule
  Uniform $\tau^{\ff}\downarrow$ &
    All edges faster free-flow &
    $\tau\downarrow\to T\downarrow\to$ access $\uparrow\to\Xi\uparrow\to$
    anchored $\tau$ partially rebounds; $L_{ij}$ shift \\
  Hub access $\tau^{\ff}\downarrow$ on hub-touching edges &
    Cheaper access to employment centres &
    Load concentrates on hub edges; $\Xi/\Xi^{\obs}$ rises on hubs;
    $\tau$ rebound can dominate local speed gain \\
  Hub/outer $\ubar,\Abar$ shocks &
    Local fundamentals &
    Minimal $\Xi$ channel; exog/endog wedges small \\
  Fringe transport &
    Outer-ring $\tau^{\ff}\downarrow$ &
    Similar to uniform but spatially localised \\
  \bottomrule
  \end{tabular}
\end{table}

\paragraph{Why exogenous vs.\ endogenous wedges differ.}
Exogenous rows isolate the \emph{first} transport effect holding $\Xi$ fixed
implicitly.
Endogenous rows add steps $\Xi\uparrow$ and $\tau$ rebound through
\eqref{eq:E6_anchored}, shrinking or reversing welfare gains when loads rise
on already-congested corridors.

\subsection{Tidal lanes}

On each treated asymmetric pair $(\text{slow},\text{fast})$:
$n_{\text{slow}}\mathrel{+}=\Delta n$, $n_{\text{fast}}\mathrel{-}=\Delta n$.

\begin{enumerate}[leftmargin=2em]
  \item \textbf{Direct mechanical}: capacity ratio $n^{\obs}/n$ changes $\Rightarrow$
        $\tau$ shifts via \eqref{eq:E6_anchored} even at fixed $\Xi$.
  \item \textbf{Network spillover}: new $\tau$ $\Rightarrow$ new $T$, re-routing,
        $\Xi$ shifts on \emph{untreated} edges too (unlike Stage~1 BPR grid).
  \item \textbf{Short-run}: $(l^R,l^F)$ pinned $\Rightarrow$ only steps 1--2;
        first-order route-cost change visible (§\ref{sec:tidal_diagnostics}).
  \item \textbf{Long-run}: sorting adjusts; sorting wedge between ST and LT.
\end{enumerate}

Typical treated-edge pattern (long-run, top-50, $\Delta n=1$):
slow direction $\Delta\tau\approx -16\%$, fast direction $\Delta\tau\approx +26\%$;
network-average $\Delta\Wbar$ remains $\lesssim 0.1\%$ because fundamentals
fixed and OD mass reweights toward hub corridors only slightly.

%====================================================================
\section{Code map}\label{sec:code}
%====================================================================

\begin{table}[htbp]
  \centering
  \caption{\texttt{data\_work/src/model/full\_estimation/} layout}
  \small
  \begin{tabular}{p{5.5cm}p{8.5cm}}
  \toprule
  Module & Role \\
  \midrule
  \texttt{build\_primitives.py} & Step~1: parquet $\to$ \texttt{Primitives} \\
  \texttt{parameters.py} & Step~2: $(\theta,\alpha,\beta,\lambda)$ \\
  \texttt{lambda\_diagnostic.py} & Step~3: cross-section $\hat\lambda$ (check only) \\
  \texttt{invert\_fundamentals.py} & Step~4: $(\ubar,\Abar,\chi)$ \\
  \texttt{equilibrium\_solve.py} & Step~5: nested loop; \texttt{congestion\_form} $\in$
        \{\texttt{legacy}, \texttt{density}, \texttt{anchored}\} \\
  \texttt{stage2\_counterfactuals.py} & Policy + tidal + structural CF orchestration \\
  \texttt{counterfactual\_levels.py} & CF1 symmetric / CF2 tidal / CF3 uniform \\
  \texttt{run\_stage2\_gate5.py} & Main driver (diagnostic-$\lambda$ or placeholder) \\
  \texttt{run\_congestion\_lambda11\_comparison.py} & Legacy vs density vs anchored \\
  \texttt{tidal\_lane\_diagnostics.py} & ST/LT congestion-cost decomposition \\
  \texttt{run\_tidal\_lane\_diagnostics\_gate5.py} & Tidal diagnostics driver \\
  \texttt{directional\_cost/} & Stage~1 static tools, \texttt{predicted\_od\_shares},
        \texttt{comparison\_summary} \\
  \bottomrule
  \end{tabular}
\end{table}

\paragraph{Main commands.}
\begin{verbatim}
# Full Stage 2 pipeline (default: diagnostic lambda two-phase)
python3 data_work/src/model/full_estimation/run_stage2_gate5.py

# Anchored lambda=1.1 tidal ST/LT diagnostics
python3 data_work/src/model/full_estimation/run_tidal_lane_diagnostics_gate5.py

# Congestion-spec comparison (legacy / density / anchored)
python3 data_work/src/model/full_estimation/run_congestion_lambda11_comparison.py
\end{verbatim}

\paragraph{Output roots.}
\begin{itemize}[leftmargin=2em]
  \item \texttt{stage2\_endogenous\_edgeproj\_v2\_lcc\_gate5\_v1/}
  \item \texttt{tidal\_lane\_diagnostics\_edgeproj\_v2\_lcc\_gate5\_v1/}
  \item \texttt{congestion\_lambda11\_comparison\_edgeproj\_v2\_lcc\_gate5\_v1/}
  \item \texttt{anchored\_model\_audit\_gate5\_v1/}
\end{itemize}

Handoff memos:
\path{Documents/L2_notes/stage2_endogenous_congestion_handoff_20260628.md},
\path{Documents/L2_notes/congestion_lambda11_density_vs_anchored_20260629.md},
\path{Documents/L2_notes/tidal_lane_st_lt_congestion_diagnostics_20260629.md}.

%====================================================================
\section{Results: baselines and policy CFs}\label{sec:results}
%====================================================================

\subsection{Baselines (anchored $\lambda_{\mathrm{phys}}=1.1$, $d=0.25$)}

\begin{table}[htbp]
  \centering
  \caption{Baselines (gate5)}
  \begin{tabular}{lrrrr}
  \toprule
  Mode & $\Wbar$ & Mean $\tau/\tau^{\obs}$ & Median & Outer iters \\
  \midrule
  Exogenous $\tau^{\obs}$ & 1.861 & 1.000 & 1.000 & 1 \\
  Endogenous anchored     & 1.861 & 1.001 & 1.000 & 36 \\
  \bottomrule
  \end{tabular}
\end{table}

Legacy unanchored + diagnostic $\lambda\approx 0.017$ (old default) pulled mean
$\tau$ $\approx 20\%$ below $\tau^{\obs}$---superseded by anchored spec.

\subsection{Policy counterfactuals (anchored endogenous vs.\ respective baseline)}

\begin{table}[htbp]
  \centering
  \caption{Policy CFs (\% change vs.\ own baseline, anchored $\lambda=1.1$)}
  \label{tab:policy}
  \small
  \begin{tabular}{lrrrr}
  \toprule
  Policy & $\Delta\Wbar$ (exog) & $\Delta\Wbar$ (endog) & OD dissim.\ (endog) &
    Mean $\Delta\Xi$ (endog) \\
  \midrule
  Uniform +10\% speed  & $+0.81$ & $+0.02$  & 0.05 & $+0.4$ \\
  Hub access +15\%     & $+0.16$ & $+0.01$  & 0.05 & $+0.1$ \\
  Hub productivity +10\% & $-0.68$ & $-0.59$ & 0.40 & $-1.4$ \\
  Outer amenity +10\%  & $+0.35$ & $+0.30$  & 0.20 & $+0.8$ \\
  Fringe transport +20\% & $+0.56$ & $+0.58$ & 0.54 & $+20$ \\
  \bottomrule
  \end{tabular}
\end{table}

Under anchored baseline, transport-policy $\Delta\Wbar$ shrink toward zero
relative to legacy unanchored runs; interpretation shifts to route-cost and
$\Xi/\Xi^{\obs}$ objects.

%====================================================================
\section{Structural and tidal counterfactuals}\label{sec:structural}
%====================================================================

\paragraph{CF1.} Symmetric corridor cost: initialise $\tau$ from geometric mean
of $(\tau^{\obs}_{kl},\tau^{\obs}_{lk})$ on reciprocal pairs; re-solve anchored
equilibrium. $\Delta\Wbar\approx 0$ (symmetrisation negligible at gate5 scale).

\paragraph{CF2--CF3.} Tidal top-$K$ and uniform lane expansion ($\times 1.10,1.25,1.50$).

\paragraph{Tidal ST/LT (anchored, $\lambda=1.1$).}

\begin{table}[htbp]
  \centering
  \caption{Tidal lanes: welfare and congestion cost}
  \label{tab:tidal}
  \small
  \begin{tabular}{lrrrrr}
  \toprule
  Spec & $\Delta\Wbar$ ST & $\Delta\Wbar$ LT &
    $\Delta$cost ST (s) & $\Delta$cost LT (s) & Pairs \\
  \midrule
  top-25, $\Delta n=1$  & $+0.04$ & $-0.03$ & $-0.12$ & $-0.44$ & 25 \\
  top-50, $\Delta n=1$  & $+0.13$ & $-0.06$ & $-0.42$ & $-0.83$ & 50 \\
  top-100, $\Delta n=1$ & $+0.28$ & $-0.12$ & $-0.72$ & $-1.24$ & 100 \\
  \bottomrule
  \end{tabular}
\end{table}

$\Delta$cost = change in OD-weighted average bilateral route time holding
baseline shares (ST) or with long-run shares (LT); see §\ref{sec:tidal_diagnostics}.

Exogenous BPR grid (Stage~1 fixed-$\xi$, no spillover): 14 specs $\times$ ST/LT;
$\Delta\Wbar\in[-0.04\%,+0.01\%]$.

%====================================================================
\section{Tidal diagnostics: ST/LT decomposition}\label{sec:tidal_diagnostics}
%====================================================================

Beyond $\Wbar$, \texttt{tidal\_lane\_diagnostics.py} reports:

\begin{align}
  \Delta C^{\mathrm{ST}} &= \sum_{ij} s^{0}_{ij}\,\big(T^{\mathrm{ST}}_{ij}-T^{0}_{ij}\big),
    \label{eq:ST_cost}\\
  \Delta C^{\mathrm{LT}} &= \sum_{ij} s^{\mathrm{LT}}_{ij}\,T^{\mathrm{LT}}_{ij}
    - \sum_{ij} s^{0}_{ij}\,T^{0}_{ij},
    \label{eq:LT_cost}
\end{align}
where $T_{ij}=\tau_{ij}$ (minutes), $s_{ij}$ model-implied OD shares.

Decomposition:
\begin{align*}
  \Delta C^{\mathrm{LT}}
  &= \underbrace{\sum_{ij}s^{0}_{ij}(T^{\mathrm{ST}}_{ij}-T^{0}_{ij})}_{\text{direct ST (pinned sorting)}}
   + \underbrace{\sum_{ij}s^{0}_{ij}(T^{\mathrm{LT}}_{ij}-T^{\mathrm{ST}}_{ij})}_{\text{GE cost wedge}}
   + \underbrace{\sum_{ij}(s^{\mathrm{LT}}_{ij}-s^{0}_{ij})T^{\mathrm{LT}}_{ij}}_{\text{sorting wedge}}.
\end{align*}

Treated-edge tables report slow/fast $\Delta\tau$ in minutes and percent.

%====================================================================
\section{Model validation and implied flows}\label{sec:validation}
%====================================================================

\subsection{Marginal vs.\ bilateral fit}

\begin{table}[htbp]
  \centering
  \caption{OD fit (anchored baseline vs.\ observed $L_{ij}$)}
  \begin{tabular}{lr}
  \toprule
  Metric & Value \\
  \midrule
  Global L1 (pred vs.\ obs) & $\approx 69\%$ \\
  Off-diagonal L1 & $\approx 51\%$ \\
  Intrazonal share: obs / pred & $8.0\%$ / $54.9\%$ \\
  Weighted log correlation & $\approx 0.59$ \\
  Residence marginal L1 & $\approx 0.02\%$ \\
  Workplace marginal L1  & $\approx 0.03\%$ \\
  \bottomrule
  \end{tabular}
\end{table}

Anchoring $\tau$ does \textbf{not} improve OD fit (same as exogenous $\tau^{\obs}$).
Primary miss: Leontief diagonal $T_{ii}$ inflates intrazonal $\widehat{L}_{ii}$.
See \path{Documents/L2_notes/stage1_od_fit_diagnosis_and_handoff_20260628.md}.

\subsection{Implied link traffic $\Xi$}

At anchored baseline, $\Xi$ matches $\Xi^{\obs}$ (median deviation $0\%$,
p95 $|\Delta|<2.4\%$).

\paragraph{Spatial pattern (qualitative).}
\begin{itemize}[leftmargin=2em]
  \item Hub-touching edges: median $\Xi\approx 49$ vs.\ outer fringe $\approx 2.3$.
  \item Top-loaded corridors cluster near high-$L^F$ grids (e.g.\ \texttt{sq\_22\_39}).
  \item $\corr(\ln\Xi,\ln L^F_{\mathrm{dest}})\approx +0.46$; $\corr(\ln\Xi,\ln\Abar_{\mathrm{dest}})\approx +0.43$.
\end{itemize}

\paragraph{Not validated.}
\begin{itemize}[leftmargin=2em]
  \item $\corr(\ln\Xi,\ln(\tau^{\obs}/\tau^{\ff}))\approx +0.08$.
  \item $\Xi$ should \textbf{not} be reported as ``traffic count'' in prose.
\end{itemize}

\subsection{Routine checks (checklist)}

\begin{enumerate}[leftmargin=2em]
  \item[\checkmark] Marginal inversion residuals.
  \item[\checkmark] $\tau$ vs.\ $\tau^{\obs}$ under anchored baseline.
  \item[\checkmark] $\Xi$ vs.\ $\Xi^{\obs}$ at baseline.
  \item[\checkmark] OD fit decomposition (global / off-diag / intra).
  \item[\checkmark] Exogenous vs.\ endogenous baseline comparison.
  \item[\checkmark] Congestion-spec robustness (legacy / density / anchored).
  \item[\checkmark] Tidal ST/LT cost decomposition.
  \item[$\circ$] Intrazonal workplace-choice specs (\texttt{exclude\_diag\_plus\_proxy}).
  \item[$\circ$] External validation (Shenzhen/Baidu corridor shocks).
  \item[$\circ$] Structural $\lambda$ from panel link data.
\end{enumerate}

Audit CSV:
\texttt{data\_work/outputs/anchored\_model\_audit\_gate5\_v1/}.

%====================================================================
\section{Interpretation for the paper}\label{sec:interpretation}
%====================================================================

\begin{enumerate}[leftmargin=2em]
  \item \textbf{Stage~2 adds feedback} \eqref{eq:feedback_loop}; Stage~1 is the
        exogenous-$\tau$ limit.
  \item \textbf{Anchored congestion} is the operational spec: baseline travel
        times match data; $\lambda=1.1$ applies to \emph{relative} capacity/load
        changes---consistent with Barwick et al.\ and exact-hat CF logic.
  \item \textbf{$\Wbar$ alone understates tidal policy}: report ST/LT route-cost
        seconds and treated-edge $\Delta\tau\%$.
  \item \textbf{Do not claim} full OD replication or that $\Xi$ matches road
        sensors; do claim marginal calibration and relative counterfactual
        transport costs.
  \item \textbf{Next estimation}: panel/IV $\lambda$; optional SMM OD moments;
        Shenzhen tidal experiment as external check on CF \emph{signs}.
\end{enumerate}

%====================================================================
\section{Next steps}\label{sec:next}
%====================================================================

\begin{enumerate}[leftmargin=2em]
  \item Wire \texttt{congestion\_form=anchored} as default in
        \texttt{run\_stage2\_gate5.py} (replace diagnostic-$\lambda$ two-phase).
  \item Figures: $\Xi$ heat map; tidal treated corridors; ST/LT cost bars.
  \item Align $\theta,\alpha,\beta$ with coauthor calibration.
  \item Directed vs.\ symmetric \emph{calibration world} under anchored congestion.
  \item Manuscript merge: cross-ref Stage~1
        \path{FixedDirectedCost_Stage1_gate5.tex} for data and Stage~1 CF baselines.
\end{enumerate}

\end{document}
